\chapter{A Short History of Randomized Computation}
\label{ch:history}

The question whether randomness helps computation is nearly as old as the question whether $\Pc = \NP$. This chapter sketches its history, from Valiant\index{Valiant, Leslie}'s counting classes through the primality tests to the hardness versus randomness programme, and the parallel history of Monte Carlo counting that leads to this book. It is not a substitute for Sipser\index{Sipser, Michael}'s survey of the $\Pc$ versus $\NP$ question, but it fixes the names and dates that the rest of the book refers to.

\section{Valiant and counting}

Leslie Valiant's two papers of 1979 opened the complexity theory of counting. The first defined $\sharpP$ and showed that the permanent of a zero one matrix is $\sharpP$ complete. The second showed that many counting problems whose decision versions are easy are nonetheless $\sharpP$ complete: counting perfect matchings in bipartite graphs, counting satisfying assignments of monotone formulas in two conjunctive normal form, counting independent sets. The phenomenon that decision can be easy while counting is hard was new and surprising. The permanent is the standard example: a bipartite graph has a perfect matching if and only if its permanent is positive, which is decided by Edmonds\index{Edmonds, Jack}'s algorithm, but the value of the permanent is as hard to compute as anything in $\sharpP$.

Toda\index{Toda's theorem}'s theorem of 1991, that the polynomial hierarchy is contained in $\Pc^{\sharpP}$, placed counting above the hierarchy. The proof goes through the class $\oplus\Pc$ of parity problems and the Valiant Vazirani theorem, which we discuss in Chapter~\ref{ch:pseudo}, and it means that a single counting oracle is at least as powerful as any fixed number of alternating quantifiers.

\section{Randomized algorithms}

The use of randomness in algorithms predates the formal theory. Monte Carlo methods go back to Ulam\index{Ulam, Stanislaw}, von Neumann\index{von Neumann, John} and Metropolis\index{Metropolis algorithm} at Los Alamos in the 1940s, and the Metropolis algorithm of 1953 is the ancestor of every Markov chain in this book. But the recognition that randomness might be a computational resource, in the sense that random algorithms could solve problems that deterministic ones could not solve efficiently, came in the 1970s.

The decisive examples were the primality tests of Solovay and Strassen in 1977 and Rabin in 1980, building on Miller's deterministic test of 1976 that depended on the generalised Riemann hypothesis. Both tests are in $\coRP$: a composite is rejected with probability at least a half by a random witness, and a prime is never rejected. For a quarter century primality was the flagship problem in $\coRP$ not known to be in $\Pc$, and it was widely cited as evidence that randomness helps. Then in 2002 Agrawal\index{Agrawal, Kayal and Saxena}, Kayal and Saxena found a deterministic polynomial time algorithm, and primality became instead the flagship example of successful derandomization.

The formal classes were defined by Gill in 1977: $\RP$, $\coRP$, $\ZPP$, $\BPP$ and $\mathrm{PP}$. Gill proved the basic inclusions, among them $\NP \subseteq \mathrm{PP}$, where $\mathrm{PP}$ is the class of problems solvable by a probabilistic machine with no gap in the acceptance probability. That $\mathrm{PP}$ is as hard as counting in a precise sense, $\Pc^{\mathrm{PP}} = \Pc^{\sharpP}$, was observed later, after Valiant defined $\sharpP$ in 1979, and is the link to $\sharpP$ that Toda exploited. Adleman\index{Adleman's theorem}'s theorem that $\RP \subseteq \Ppoly$ appeared in 1978; Sipser's theorem that $\BPP$ is in the polynomial hierarchy, with G\'acs's improvement to the second level and Lautemann's simplified proof, appeared in 1983.

The second flagship problem, polynomial identity testing, has a longer and more tangled history. The Schwartz Zippel lemma, that a nonzero polynomial of degree $d$ over a field vanishes at a random point of a set $S^n$ with probability at most $d / |S|$, was published by Zippel and by Schwartz in 1979 and 1980, and had been observed earlier by DeMillo and Lipton. Its consequence, that the identity of two arithmetic circuits can be tested by evaluating at a random point, is in $\coRP$. It remains outside $\Pc$ as far as anyone knows, and Kabanets and Impagliazzo\index{Kabanets and Impagliazzo} showed in 2004 why: a deterministic algorithm would prove circuit lower bounds that nobody can prove.

\section{Hardness versus randomness}

The idea that randomness can be removed at the price of a hardness assumption is due to Yao in 1982, who showed that a one way permutation yields a pseudorandom generator in the cryptographic sense, and to Blum and Micali in the same year. Nisan and Wigderson\index{Nisan and Wigderson} in 1994 found that for derandomization, as opposed to cryptography, a much weaker assumption suffices: it is enough that some function in exponential time be hard for circuits of a given size, and the generator need only fool circuits of that size. Their generator stretches a seed of length $O(\log n)$ to $n$ bits by evaluating the hard function on nearly disjoint subsets of the seed, and the analysis shows that a distinguisher for the output would yield a small circuit for the hard function.

Impagliazzo and Wigderson in 1997 closed the gap. By derandomizing Yao's XOR lemma, which amplifies mild hardness to extreme hardness, they showed that if some language in $\mathrm{E}$ requires circuits of size $2^{\Omega(n)}$ then $\Pc = \BPP$. The hypothesis is a strong form of the belief that exponential time is really exponential. It is believed. It is not proved, and proving it would be a major advance in circuit complexity, a field in which the best known lower bounds for explicit functions are only linear.

Impagliazzo's survey of 1995 named five possible worlds according to which of the natural hardness assumptions hold: Algorithmica, in which $\Pc = \NP$; Heuristica, in which $\NP$ problems are hard in the worst case but easy on average; Pessiland, in which hard on average problems exist but one way functions do not; Minicrypt, in which one way functions exist but public key cryptography does not; and Cryptomania, in which everything exists. In every one of the five worlds, $\Pc = \BPP$ is either proved or widely expected. There is no named world in which randomness is a genuine resource, which is a measure of how strongly the field believes that it is not.

\section{Monte Carlo counting}

The thread that leads to this book runs in parallel with the complexity thread and intersects it at several points. The method of counting by sampling is due to Knuth in 1975, who estimated the size of a backtracking tree by random probing, and to Kahn and Harris in 1951, who estimated rare event probabilities in particle transport by splitting trajectories. Jerrum, Valiant and Vazirani in 1986 made the connection between counting and sampling precise and placed it in the complexity framework. Jerrum and Sinclair in 1989, Dyer, Frieze and Kannan in 1991, and Jerrum, Sinclair and Vigoda in 2004 supplied the rapid mixing bounds that turned the connection into algorithms. Diaconis\index{Diaconis and Holmes} and Holmes in 1995 wrote the paper that first treated the queens as a Monte Carlo counting problem and that we took as our starting point.

On the statistical side, the sequential methods of annealed importance sampling, nested sampling\index{nested sampling}, and the various energy level samplers of physics all compute normalising constants by walking through a ladder of distributions, and Polson and Scott in 2014 observed that the ladder is always a traversal of the quantile function of the likelihood ordinate. The present book is the result of noticing that the quantile function is also the object that the complexity theorists are looking at when they ask whether a randomized algorithm can be derandomized.
